% tr-05 片段 c：原文 pp.41–43（PDF pp.42–44）
% 内容：5.2 超势的产生（步骤 1–4，(5.21)–(5.26)）→ 5.3 节标题 → (5.27) 止
% [原文 p.41]

\subsection{超势的产生}

对偶系统的拉氏量 (5.17)–(5.19) 当然不是精确的。在对偶化的过程中我们忽略了许多东西。例如，当我们在 (5.2) 中积掉 $J$-场时，我们假装 $\rho$ 是一个固定的非零常数，但它实际上是一个涨落的场。纳入修正的通常做法是：把场在其经典值附近展开（例如在 $\rho_i^2=\Re(y_i)$ 取大值处展开），然后对展开的每一项作路径积分。这类修正称为\emph{微扰修正}，它们取“耦合常数”的幂级数形式，在目前情形下该耦合常数就是 $1/\Re(y_i)$。对 Kähler 势 (5.18) 来说，这样的修正会多得不可胜数。然而对扭超势 (5.19) 来说，微扰修正没有任何余地——根本不可能写出一个与全纯性、（可能具有异常的）R-对称性以及周期性 $y_i\equiv y_i+2\pi\ii$ 相容的修正项。也就是说，(5.19) 在微扰意义下是精确的。

然而，扭超势可以有\emph{非微扰}修正。典型的非微扰效应由瞬子产生。在像我们的 GLSM 这样规范对称性被破缺的二维阿贝尔规范理论中，瞬子就是\emph{涡旋}。涡旋是带荷标量 $\phi$ 与规范势 $v$ 使作用量极小的一种位形，具有如下特征：$\phi$ 在某点 $p$ 处为零；在 $p$ 周围一个小圆盘 $D$ 之外，$|\phi|$ 几乎等于（非零的）真空值，但 $\phi$ 的相位沿着绕 $D$ 的圆周有非平凡的绕数；规范势在 $D$ 之外几乎平坦（$\dd v=0$），且在那里 $\phi$ 关于 $v$ 几乎平行。单个涡旋可以看作“为 $\phi$ 的相位 $\varphi$ 创造一单位绕数”的东西。既然 T-对偶交换绕数与动量，它就应对偶于“创造一单位动量”的东西，而后者由 $\ee^{\pm\ii\tphi}$ 提供。A-手征算子 $\ee^{\pm y}$ 是它的超对称完备化，而其中之一

\begin{equation}
    \ee^{-y},
    \tag{5.21}
\end{equation}
具有矢量 R-荷 $0$ 与轴 R-荷 $2$（见 (5.15)），因此可以加到扭超势中。我们将证明：这种形式的修正确实由非微扰效应产生。事实上，我们将证明该修正就是这类项之和（每个 $\phi_i$ 对应一项），而精确的扭超势为

\begin{equation}
    \tW=\sum_{a=1}^{k}\Bigl(\sum_{i=1}^{N}Q_i^a y_i-t^a\Bigr)\sigma_a+\mu\sum_{i=1}^{N}\ee^{-y_i}.
    \tag{5.22}
\end{equation}

% [Tutor Note: p.41 (5.21) 的 $\ee^{-y}$ 之所以能进扭超势，可用“荷 + 量纲”双重判据核对：由 (5.15) 的 $y\to y-2\ii\beta$ 得 $\ee^{-y}\to\ee^{-y}\ee^{2\ii\beta}$，即轴 R-荷 $2$，与 $\tW$ 必须具有的轴 R-荷 $2$ 相符（$\sigma$ 本身轴 R-荷 $2$、$\ee^{-\ty}$ 轴 R-荷 $0$，见步骤 2 中那句 “$\sigma$, $\ee^{-\ty}$ and $\Lambda$ has axial R-charge 2, 0 and 2”）；量纲上 $\mu$ 是一次幂，正好补足 $\ee^{-y_i}$ 的无量纲性。原文 p.41 “In a two-diemsional Abelian gauge theory” 与 “with repect to $v$” 均为排印笔误（应作 two-dimensional / respect），正文照译。]

% [原文 p.42]

\para{步骤 1——归约到 $\CP^0$ 模型}

首先我们证明：若“(5.22) 是精确扭超势”这一断言对 $\CP^0$ 模型成立，则它对一般模型也必定成立。我们从 $N$ 个 $\CP^0$ 模型的直积出发。在假定该断言成立的前提下，对其中每一个施加 T-对偶，便得到 $N$ 份对偶理论的直积。特别地，我们有 $\tilde{K}=\sum_{i=1}^{N}\tilde{K}_i$ 以及 $W=\sum_{i=1}^{N}W_i$，其中 $\tilde{K}_i$ 是 $\sigma_i$ 与 $y_i$ 的某个复杂函数，而 $W_i=(y_i-t_i)\sigma_i+\mu\,\ee^{-y_i}$。接着，我们通过改变规范多重态的动能项来形变原理论，即形变 $K_{\mathrm{gauge}}=\sum_i|\sigma_i|^2/(2e_i^2)$。为此我们取分解 $U(1)^N\cong G\times U(1)^{N-k}$，并相应地写 $\sigma_i=\sum_{a=1}^{k}Q_i^a\sigma_a+\sum_{f=k+1}^{N}P_i^f\sigma_f$。然后把 $K_{\mathrm{gauge}}$ 形变为

\begin{equation}
    K'_{\mathrm{gauge}}=\sum_{a=1}^{k}\frac{1}{2e_a^2}|\sigma_a|^2+\sum_{f=k+1}^{N}\frac{1}{2e_f^2}|\sigma_f|^2.
    \tag{5.23}
\end{equation}

这一形变当然会改变 $\tilde{K}$，但不可能改变 $\tW$，因为形变参数不是 A-手征的。也就是说，$\tW$ 仍然是那个和 $\sum_{i=1}^{N}\tilde{W}_i$。在这个阶段，我们取 $e_f\to0$（$f=k+1,\dots,N$）的极限。这就冻结了对应于 $U(1)^{N-k}$ 的规范多重态分量，原系统便约化为所考虑规范群为 $G$ 的系统。标度 $\mu$ 处的 FI-theta 参数是 $t^a=\sum_{i=1}^{N}Q_i^a t_i$，这可以由匹配有效扭超势看出。在对偶一侧，$\tW$ 保持不变，只是 $\sigma_i$ 现在受到约束 $\sigma_i=\sum_{a=1}^{k}Q_i^a\sigma_a$，而这正是 (5.19) 中的那个。

% [Tutor Note: p.42 步骤 1 中 “we have $\tilde{K}=\sum_i\tilde{K}_i$ and $W=\sum_i W_i$ ... and $W_i=(y_i-t_i)\sigma_i+\mu\ee^{-y_i}$” 里的 $W,W_i$ 原文就漏写了波浪号（已用 900 dpi 裁图逐字核对：该处 $K$ 带 $\sim$ 而 $W$ 不带）；按上下文（同一段下一句 “$\tW$ remains to be the sum $\sum_{i=1}^{N}\tilde{W}_i$”）它们应是 $\tW,\tW_i$。译文照原文保留 $W$，此处指出。同页 “it must be ture for the general model” 应作 true、“This defomation” 应作 deformation，均为排印笔误。]

\para{步骤 2——排除其他修正}

其次我们证明：在 $\CP^0$ 模型中，

\begin{equation}
    \Delta\tW=C\cdot\mu\,\ee^{-y}
    \tag{5.24}
\end{equation}
（其中 $C$ 为某个常数）是扭超势唯一可能的修正。扭超势 $\tW$ 必须是 RG 不变的 A-手征算子与参数的函数，且必须全纯；这些算子与参数是 $\sigma$、$\ty=y-t$ 与 $\Lambda=\ee^{-t}\mu$。既然 $\ty$ 具有周期性 $\ty\equiv\ty+2\pi\ii$，它就必须通过 $\ee^{-\ty}$ 依赖于 $\ty$，唯一例外是经典项 $\ty\sigma$——后者是可容许的，因为 theta 参数平移 $2\pi$ 是可容许的。此外，它必须具有轴 R-荷 $2$，而这是以 FI-theta 参数按 $t\to t-2\ii\beta$ 变换为前提的。注意 $\sigma$、$\ee^{-\ty}$ 与 $\Lambda$ 分别具有轴 R-荷 $2$、$0$ 与 $2$。由这些条件，可能的修正 $\Delta\tW$ 必须取 $\Lambda f(\sigma/\Lambda,\ee^{-\ty})$ 的形式，其中 $f$ 是某个全纯函数。
% [原文 p.43]
由于规范对称性被 $\phi\neq0$ 破坏，扭超势在 $\sigma=0$ 处必须是正则的。再者，当 $\Re(y)$ 取大值（对应于 $|\phi|$ 取大值）时修正应当很小。因此修正不可能在 $\Re(\ty)\to+\infty$ 时增长。于是函数 $f(z,w)$ 在 $(z,w)=(0,0)$ 处必须解析，并可以按 $z$ 与 $w$ 的幂级数展开：

\begin{equation}
    \Delta\tW=\Lambda\sum_{n,m\ge0}C_{n,m}(\sigma/\Lambda)^n\ee^{-m\ty}
        =\sum_{n,m\ge0}C_{n,m}\sigma^n\mu^{1-n}\ee^{-t+(n+m)t}\ee^{-my}.
    \tag{5.25}
\end{equation}
此外，当 $\zeta=\Re(t)$ 取大值时修正应当很小，因此只有 $n+m\le1$ 的项能被生成。$(n,m)=(0,0)$ 是常数，可以忽略。$(n,m)=(1,0)$ 与经典项形式相同，故必须为零。于是我们只能有 $(n,m)=(0,1)$，也就是 (5.24)。

% [Tutor Note: p.42–43 式 (5.25) 的第二行可逐步核对：$\Lambda(\sigma/\Lambda)^n=\Lambda^{1-n}\sigma^n=(\ee^{-t}\mu)^{1-n}\sigma^n=\sigma^n\mu^{1-n}\ee^{-t(1-n)}$，而 $\ee^{-m\ty}=\ee^{-m(y-t)}=\ee^{-my}\ee^{mt}$；两者相乘即得 $\sigma^n\mu^{1-n}\ee^{-t+(n+m)t}\ee^{-my}$，其中 $-t(1-n)=t(n-1)$ 与 $+mt$ 合并成 $-t+(n+m)t$。同一核对也说明为何只保留 $n+m\le1$：$\zeta=\Re(t)$ 大时该组合给出 $\ee^{(n+m-1)\zeta}$，只有 $n+m\le1$ 的项不被压低。原文 p.43 “Furtheremore” 应作 Furthermore、“genertated” 应作 generated（排印笔误）。]

\para{步骤 3——瞬子演算}

现在我们来证明 (5.24) 确实以一个非零系数 $C$ 被生成。

\para{步骤 4——系数}

最后我们证明系数 $C$ 等于 $1$。让我们在 $\sigma$ 取大值处积掉 $y$ 场。这通过解方程 $\partial_y\tW=0$ 求出 $y$、再把答案代回 $\tW$ 来完成。该方程是 $\sigma-C\mu\,\ee^{-y}=0$，解为 $y=-\log(\sigma/C\mu)$。把它代回 $\tW$，得到

\begin{equation}
    \tW=\bigl(-\log(\sigma/C\mu)-t\bigr)\sigma+\sigma=-t\sigma-\sigma\bigl(\log(\sigma/C\mu)-1\bigr).
    \tag{5.26}
\end{equation}

这正是 $\CP^0$ 模型在大 $\sigma$ 处的有效扭超势，见 (4.16)——前提是 $C$ 为 $1$。这就证明了 $C=1$。我们也可以从 (5.22) 出发，在一般模型中得到正确的有效扭超势 (4.16)。

% [Tutor Note: p.43 式 (5.26) 中那个单独的 “$+\sigma$” 不是笔误：把 $y=-\log(\sigma/C\mu)$ 代回 $\tW=(y-t)\sigma+C\mu\ee^{-y}$ 时，最后一项给出 $C\mu\cdot(\sigma/C\mu)=\sigma$，与第一项合并即得 (5.26) 的第二个等号。核对判据：若漏掉这个 $+\sigma$，得到的是 $-t\sigma-\sigma\log(\sigma/C\mu)$，与 (4.16) 的 $\CP^0$ 有效扭超势（含常数项、形如 $-\sigma(\log(\sigma/\mu)-1)$）不符；正是 $+\sigma$ 提供了那个常数项 $1$。]

\subsection{镜像对称的推导一——无超势的理论}

让我们积掉 $U(1)^k$ 规范多重态。这给出约束

\begin{equation}
    \sum_{i=1}^{N}Q_i^a y_i-t^a=0,\qquad a=1,\dots,k,
    \tag{5.27}
\end{equation}
